\section{Contexte et analyse des besoins}
\markboth{1. Contexte et analyse des besoins}{}

\subsection{Présentation du contexte}

\subsubsection*{L'organisation et son obligation réglementaire}

La Ville de Seattle impose depuis 2010, par ordonnance municipale, que tout
bâtiment non résidentiel ou résidentiel collectif de plus de
\num{20000}~pieds carrés déclare annuellement sa consommation énergétique. Ce
dispositif, le \emph{Building Energy Benchmarking}, sert un objectif public
affiché : atteindre la neutralité carbone à l'échelle de la ville. Les données
collectées sont publiées en open data, millésime par millésime.

Le service porteur — l'\emph{Office of Sustainability \& Environment} — dispose
donc d'un inventaire déclaratif annuel de plusieurs milliers de bâtiments, avec
leurs caractéristiques structurelles et leurs consommations relevées.

\subsubsection*{Maturité data et point de départ}

La maturité data de l'organisation est asymétrique, et c'est ce déséquilibre qui
fonde le besoin.

\begin{itemize}[leftmargin=1.4em,itemsep=0.25em]
  \item \textbf{Forces} : une collecte structurée, annuelle, avec un référentiel
        stable de typologies de bâtiments ; une publication en open data qui
        rend les données auditables.
  \item \textbf{Faiblesses} : la donnée est \emph{publiée} mais peu
        \emph{exploitée}. Elle décrit le passé des bâtiments déclarants, et ne
        dit rien de ceux qui ne déclarent pas — ni de ce qu'il faudrait
        prioriser.
\end{itemize}

\begin{constat}{Le point de départ technique}
Un projet antérieur de modélisation a produit deux modèles de régression sous
forme de notebooks Jupyter : consommation d'énergie et émissions de gaz à effet
de serre, à partir des caractéristiques structurelles. Ces notebooks constituent
\textbf{la solution existante auditée en section~2}. Ils n'ont jamais été
déployés, ni suivis, ni testés.
\end{constat}

\subsubsection*{Périmètre de ce rapport}

Ce projet est un projet personnel construit sur des données publiques. Le
contexte organisationnel décrit ci-dessus est réel — l'ordonnance, le programme
de benchmarking et les données le sont — mais les parties prenantes n'ont pas été
interrogées : leurs besoins ont été \textbf{reconstitués} à partir de la
documentation publique du programme et de la nature même des données publiées.
Cette limite est assumée et signalée ici plutôt que masquée ; elle est rappelée
en section~\ref{sec:limites-demarche}.

\subsection{Collecte et analyse du besoin métier}

\subsubsection*{Parties prenantes identifiées}

\begin{table}[H]
\centering
\small
\begin{tabularx}{\linewidth}{@{}L{3.6cm}XL{4.2cm}@{}}
\toprule
\textbf{Partie prenante} & \textbf{Attente vis-à-vis de la donnée} & \textbf{Contrainte propre} \\
\midrule
Service développement durable & Cibler les bâtiments à auditer en priorité, avec des moyens d'audit finis & Budget d'audit contraint, obligation de justifier le ciblage \\
Propriétaires et gestionnaires & Situer leur bâtiment par rapport à des bâtiments comparables & Ne disposent pas toujours de relevés exploitables \\
Bureaux d'études énergétiques & Pré-qualifier un parc avant intervention & Ont besoin d'un ordre de grandeur, pas d'une valeur contractuelle \\
Citoyens et élus & Transparence sur l'empreinte du parc & Données publiques, donc opposables \\
\bottomrule
\end{tabularx}
\caption{Parties prenantes et attentes reconstituées à partir de la documentation du programme.}
\end{table}

\subsubsection*{Hiérarchisation des besoins}

Les besoins ont été classés selon deux axes : la valeur pour le service porteur
et la faisabilité au regard des données réellement disponibles. Ce croisement
élimine d'emblée les besoins que la donnée ne peut pas servir, quelle que soit
leur valeur.

\begin{table}[H]
\centering
\small
\begin{tabularx}{\linewidth}{@{}L{5.4cm}C{2.2cm}C{2.4cm}L{3.6cm}@{}}
\toprule
\textbf{Besoin} & \textbf{Valeur} & \textbf{Faisabilité} & \textbf{Décision} \\
\midrule
Estimer la consommation d'un bâtiment non relevé & Élevée & Élevée & \textbf{Retenu — cœur du projet} \\
Prioriser les audits sur un parc entier & Élevée & Élevée & \textbf{Retenu} \\
Quantifier la fiabilité de chaque estimation & Élevée & Moyenne & \textbf{Retenu — différenciant} \\
Expliquer ce qui porte une estimation & Moyenne & Élevée & Retenu \\
Prédire l'effet d'une rénovation & Élevée & \textbf{Nulle} & Écarté : aucune variable de rénovation dans les données \\
Suivre la dérive du parc dans le temps & Moyenne & Faible & Écarté : un millésime annuel, un seul utilisateur \\
\bottomrule
\end{tabularx}
\caption{Matrice de priorisation des besoins.}
\end{table}

\begin{constat}{Le besoin qui a réorienté tout le projet}
Estimer la consommation d'un bâtiment \textbf{non relevé} n'est pas une variante
du problème initial : c'est une contrainte qui interdit d'utiliser toute variable
dérivée de la consommation. Or la solution existante en utilisait. Ce constat est
l'origine de l'écart le plus structurant du projet (section~\ref{sec:fuite}).
\end{constat}

\subsubsection*{Contraintes}

\begin{description}[leftmargin=0em,style=nextline,itemsep=0.35em]
  \item[Métier] L'outil doit servir la priorisation, pas la facturation. Une
        estimation à un facteur~2 près est exploitable pour classer, inutilisable
        pour contractualiser. Cette frontière doit être écrite dans l'interface,
        pas seulement dans la documentation.
  \item[Données] Un seul millésime exploité (2016), \num{1655} bâtiments après
        nettoyage, aucune variable d'usage réel (occupation, horaires,
        équipements). Le réentraînement est annuel, au rythme de publication.
  \item[Opérationnelle] Un seul utilisateur métier, pas d'équipe d'exploitation.
        Toute brique de supervision ajoutée devra être maintenue par cette même
        personne — ce qui pèse dans les arbitrages de la section~\ref{sec:monitoring}.
  \item[Réglementaire] Les données sont publiques et les estimations sont
        opposables. Un chiffre affiché sans sa marge d'erreur serait, dans ce
        contexte, une information trompeuse.
\end{description}

\newpage
